\section{Network Architecture}


\subsection{The Three Core Chains}

Lux is organized as a multi-chain platform. The three core chains are
the Platform Chain (P-Chain), the Quantum Chain (Q-Chain), and the
Zero-Knowledge Chain (Z-Chain).

\begin{description}[nosep,leftmargin=0pt]
  \item[P-Chain (Platform)] Manages the validator set, staking,
        rewards, and chain creation. Uses BLS12-381 aggregate
        signatures for block finality. This is the coordination layer
        of the network.
  \item[Q-Chain (Quantum)] Manages post-quantum threshold signatures
        via Pulsar (Ring-LWE threshold sigs~\cite{corona}) and
        ML-DSA~\cite{mldsa}. Q-Chain is where the network's
        quantum-resistant cryptographic state lives.
  \item[Z-Chain (Zero-Knowledge)] Manages shielded transactions and
        privacy-preserving smart contracts through zero-knowledge
        proofs. Shielded payments, private security token transfers,
        and compliant confidential orders live here.
\end{description}

Together, P, Q, and Z form the \textbf{PQZ stack}: the Platform chain
coordinates, the Quantum chain provides post-quantum cryptographic
finality, and the Z chain provides privacy. Block finality combines
attestations from all three: a block is final when it carries a
valid BLS aggregate (P-Chain), a valid Pulsar threshold signature
(Q-Chain), and a valid ZK proof of consistency with prior state
(Z-Chain).

\subsection{Application Chains}

Application developers deploy sovereign L1 chains (\emph{subnets} in
earlier terminology, \emph{track-chains} in the current vocabulary)
via the P-Chain. Each chain chooses its own VM, validator set,
genesis parameters, and regulatory configuration. A bank issuing
tokenized treasuries runs a different chain from a gaming studio
issuing NFTs; they share the P-Chain's security and the Q/Z-Chains'
cryptographic primitives without interfering with each other.

At launch, Lux supports the following VM types:

\begin{itemize}[nosep]
  \item \textbf{EVM} (C-Chain): Ethereum-compatible execution with
        additional precompiles for PQ signatures, FHE, and threshold
        verification.
  \item \textbf{AVM} (X-Chain): UTXO-based asset exchange chain.
        Inherently quantum-resistant through address hashing.
  \item \textbf{Custom}: Any application developer can define their
        own VM via the Subnet VM interface.
\end{itemize}

\subsection{Consensus}

Lux uses a family of related consensus protocols---collectively known
as \emph{Lux Consensus}---that share a common structure: probabilistic
metastable voting, DAG-based transaction ordering, and leaderless
finality. The protocol family is named after physical phenomena:
\emph{Photon}, \emph{Wave}, \emph{Flare}, \emph{Focus}, \emph{Horizon},
\emph{Ray}, \emph{Nova}, \emph{Nebula}, \emph{Prism}, and
\emph{Quasar}. Quasar~\cite{lp105} is the composite protocol that
binds P-Chain BLS signatures, Q-Chain Corona/ML-DSA threshold
signatures, and Z-Chain ZK proofs into a single triple-proof finality
guarantee.

The consensus protocol achieves sub-second finality with no leader
election. A block is considered final when it has accumulated
attestations from at least $\frac{2n}{3} + 1$ validators across all
three core chains. The mathematical treatment is developed in
LP-020~\cite{lp105} and the accompanying protocol papers.

%% ============================================================
